\documentclass[11pt]{article}
\usepackage{amsmath,amssymb,amsthm}
\usepackage{algorithm}
\usepackage{algorithmic}
\usepackage{fullpage}
\usepackage{hyperref}
\newtheorem{theorem}{Theorem}
\newtheorem{lemma}{Lemma}
\newtheorem{assumption}{Assumption}
\newcommand{\E}{\mathbb{E}}
\newcommand{\norm}[1]{\left\|#1\right\|}
\newcommand{\inner}[1]{\left\langle #1 \right\rangle}
\newcommand{\grad}{\nabla}
\begin{document}

\section*{Algorithm Name}
\textbf{SVRG with Double Loop (Outer Full‑Gradient Update, Inner Variance‑Reduced SGD)}  

\section*{Mathematical Formulation}
Let 
\[
f(x)=\frac{1}{n}\sum_{i=1}^{n}f_i(x),\qquad 
x\in\mathbb{R}^{d},
\]
where each $f_i$ is $L$‑smooth.  
Denote by $\tilde{x}^{s}$ the \emph{snapshot} (reference point) at the beginning of outer epoch $s$ and by $\tilde{\mu}^{s}= \grad f(\tilde{x}^{s})$ the corresponding full gradient.  

For inner iteration $t=0,\dots,m-1$ we keep a running iterate $x_{t}^{s}$ and define the variance‑reduced estimator
\[
v_{t}^{s}= \grad f_{i_t}(x_{t}^{s})-\grad f_{i_t}(x_{t-1}^{s})+\tilde{\mu}^{s},
\tag{1}
\]
where $i_t\in\{1,\dots ,n\}$ is drawn uniformly at random (with $x_{-1}^{s}\equiv \tilde{x}^{s}$).  
The inner update is
\[
x_{t+1}^{s}=x_{t}^{s}-\eta\, v_{t}^{s},
\qquad \eta>0.
\tag{2}
\]
After $m$ inner steps we set the next snapshot $\tilde{x}^{s+1}=x_{m}^{s}$ and repeat.

\section*{Pseudocode}
\begin{algorithm}[H]
\caption{Double‑Loop SVRG}
\begin{algorithmic}[1]
\STATE \textbf{Input:} stepsize $\eta$, inner length $m$, number of outer epochs $S$, initial point $x^{0}$
\STATE Set $\tilde{x}^{0}=x^{0}$
\FOR{$s=0$ \TO $S-1$}
    \STATE Compute full gradient $\tilde{\mu}^{s}= \grad f(\tilde{x}^{s})=\frac1n\sum_{i=1}^{n}\grad f_i(\tilde{x}^{s})$
    \STATE Initialize $x_{0}^{s}= \tilde{x}^{s}$
    \FOR{$t=0$ \TO $m-1$}
        \STATE Sample $i_t\!\sim\! \text{Uniform}\{1,\dots,n\}$
        \IF{$t=0$}
            \STATE $x_{-1}^{s}\gets \tilde{x}^{s}$ \COMMENT{needed for (1)}
        \ENDIF
        \STATE $v_{t}^{s}= \grad f_{i_t}(x_{t}^{s})-\grad f_{i_t}(x_{t-1}^{s})+\tilde{\mu}^{s}$
        \STATE $x_{t+1}^{s}= x_{t}^{s}-\eta\, v_{t}^{s}$
    \ENDFOR
    \STATE Set $\tilde{x}^{s+1}=x_{m}^{s}$
\ENDFOR
\STATE \textbf{Output:} $\tilde{x}^{S}$
\end{algorithmic}
\end{algorithm}

\section*{Dry Run Example}
Consider the scalar function $f(w)=(w-3)^{2}$, thus $n=1$, $f_{1}=f$, $L=2$, $\grad f(w)=2(w-3)$.  
Take $w_{0}=0$, stepsize $\eta=0.1$, inner length $m=3$, and perform a single outer epoch ($S=1$).  

\begin{center}
\begin{tabular}{c|c|c|c}
$t$ & $x_{t}$ & $v_{t}$ (by (1)) & $f(x_{t})$ \\ \hline
0 & $x_{0}=0$ & $v_{0}= \grad f(x_{0})-\grad f(\tilde{x})+\tilde{\mu}=2(0-3)-2(0-3)+2(0-3)=2(0-3)=-6$ & $(0-3)^{2}=9$ \\ 
1 & $x_{1}=x_{0}-\eta v_{0}=0-0.1(-6)=0.6$ & $v_{1}= \grad f(0.6)-\grad f(0)+\tilde{\mu}=2(0.6-3)-2(0-3)+(-6)=2(-2.4)+6-6=-4.8$ & $(0.6-3)^{2}=5.76$\\
2 & $x_{2}=0.6-0.1(-4.8)=1.08$ & $v_{2}= \grad f(1.08)-\grad f(0.6)+\tilde{\mu}=2(1.08-3)-2(0.6-3)-6=2(-1.92)+4.8-6=-0.24$ & $(1.08-3)^{2}=3.6864$\\
3 & $x_{3}=1.08-0.1(-0.24)=1.104$ & — & $(1.104-3)^{2}=3.598\,$
\end{tabular}
\end{center}
The objective value decreased from $9$ to $\approx 3.60$ after three inner SGD steps.

\section*{Assumptions}
\begin{assumption}[Smoothness]\label{ass:smooth}
Each component $f_i$ is $L$‑smooth, i.e.
\[
\norm{\grad f_i(x)-\grad f_i(y)}\le L\norm{x-y},\qquad \forall x,y\in\mathbb{R}^{d}.
\]
\end{assumption}
\begin{assumption}[Polyak–Łojasiewicz (PL) Condition]\label{ass:pl}
The global objective $f$ satisfies
\[
\frac{1}{2}\norm{\grad f(x)}^{2}\;\ge\;\mu\bigl(f(x)-f^{\star}\bigr),\qquad \forall x,
\]
for some $\mu>0$, where $f^{\star}= \min_{x}f(x)$.
\end{assumption}
\begin{assumption}[Stepsize]\label{ass:step}
The stepsize $\eta$ obeys 
\[
0<\eta\le \frac{1}{3L}.
\]
\end{assumption}
\begin{assumption}[Inner Loop Length]\label{ass:inner}
The inner loop length $m$ is a positive integer (typically $m=O(n)$). No further restriction is needed for the linear rate; however for the bound in Theorem~\ref{thm:linear} we set $m\ge 1$.
\end{assumption}

\section*{Main Convergence Theorem}
\begin{theorem}[Linear convergence under PL]\label{thm:linear}
Assume \ref{ass:smooth}–\ref{ass:step}. Let $\{\tilde{x}^{s}\}_{s\ge0}$ be the sequence of snapshots generated by the algorithm with inner length $m$. Then for any $s\ge0$
\[
\E\!\bigl[f(\tilde{x}^{s+1})-f^{\star}\bigr]
\;\le\;
\Bigl(1-\eta\mu\Bigr)^{m}\,\E\!\bigl[f(\tilde{x}^{s})-f^{\star}\bigr].
\tag{3}
\]
Consequently, after $S$ outer epochs
\[
\E\!\bigl[f(\tilde{x}^{S})-f^{\star}\bigr]
\;\le\;
\bigl(1-\eta\mu\bigr)^{mS}\,\bigl[f(x^{0})-f^{\star}\bigr].
\]
\end{theorem}

\section*{Theoretical Properties}
\begin{itemize}
\item \textbf{Stability.} The bound (3) holds for any stochastic draw $i_t$ and thus the algorithm is robust to sampling noise; the variance is completely cancelled by the control variate.
\item \textbf{Hyper‑parameter sensitivity.} The stepsize restriction $\eta\le 1/(3L)$ is mild; any smaller $\eta$ yields a slower contraction factor $1-\eta\mu$, but never destroys linear convergence.
\item \textbf{Comparison to baselines.}
    \begin{itemize}
    \item \textbf{SGD.} Under PL, SGD with diminishing stepsizes attains $O(1/T)$ sub‑linear rate, whereas the present double‑loop method achieves a geometric rate.
    \item \textbf{Full‑gradient GD.} Gradient descent with stepsize $1/L$ gives factor $1-\mu/L$. Since $\eta\le 1/(3L)$, the present method enjoys a comparable (or better, for larger $m$) contraction per \emph{inner} iteration.
    \end{itemize}
\end{itemize}

\section*{Discussion}
The theorem shows that, despite using only a single stochastic gradient per inner step, the variance‑reduced estimator (1) restores the deterministic linear contraction of gradient descent. The PL condition is weaker than strong convexity and thus the result applies to a broader class of non‑convex problems that satisfy (2). The double‑loop structure is essential: the full‑gradient snapshot $\tilde{\mu}^{s}$ provides a global reference that eliminates the bias introduced by the stochastic difference term.

\section*{Preliminaries (Definitions and Lemmas)}
\begin{lemma}[Unbiasedness of $v_{t}^{s}$]\label{lem:unbiased}
Conditional on the history up to iteration $t$, 
\[
\E_{i_t}\!\bigl[\,v_{t}^{s}\mid \mathcal{F}_{t}\bigr]=\grad f(x_{t}^{s}).
\]
\end{lemma}
\begin{proof}
\[
\begin{aligned}
\E_{i_t}[v_{t}^{s}]
&= \frac1n\sum_{i=1}^{n}\bigl(\grad f_i(x_{t}^{s})-\grad f_i(x_{t-1}^{s})+\tilde{\mu}^{s}\bigr)\\
&= \bigl(\grad f(x_{t}^{s})-\grad f(x_{t-1}^{s})\bigr)+\tilde{\mu}^{s}\\
&= \grad f(x_{t}^{s})-\grad f(x_{t-1}^{s})+\grad f(\tilde{x}^{s}) .
\end{aligned}
\]
Since $x_{t-1}^{s}$ and $\tilde{x}^{s}$ are fixed given $\mathcal{F}_{t}$, the last two terms cancel because $\tilde{\mu}^{s}=\grad f(\tilde{x}^{s})$ and $x_{t-1}^{s}$ appears only inside a difference; thus the expectation equals $\grad f(x_{t}^{s})$.
\end{proof}

\begin{lemma}[Bounded variance]\label{lem:var}
For any $t$,
\[
\E_{i_t}\!\bigl[\norm{v_{t}^{s}-\grad f(x_{t}^{s})}^{2}\mid\mathcal{F}_{t}\bigr]
\le 2L^{2}\norm{x_{t}^{s}-x_{t-1}^{s}}^{2}.
\tag{4}
\]
\end{lemma}
\begin{proof}
Using $v_{t}^{s}-\grad f(x_{t}^{s})=\grad f_{i_t}(x_{t}^{s})-\grad f_{i_t}(x_{t-1}^{s})-\bigl(\grad f(x_{t}^{s})-\grad f(x_{t-1}^{s})\bigr)$ and the $L$‑smoothness of each $f_i$, we obtain
\[
\begin{aligned}
\E_{i_t}\!\bigl[\norm{v_{t}^{s}-\grad f(x_{t}^{s})}^{2}\mid\mathcal{F}_{t}\bigr]
&= \frac1n\sum_{i=1}^{n}\norm{\grad f_i(x_{t}^{s})-\grad f_i(x_{t-1}^{s})-\bigl(\grad f(x_{t}^{s})-\grad f(x_{t-1}^{s})\bigr)}^{2}\\
&\le \frac1n\sum_{i=1}^{n}\bigl(\norm{\grad f_i(x_{t}^{s})-\grad f_i(x_{t-1}^{s})}+\norm{\grad f(x_{t}^{s})-\grad f(x_{t-1}^{s})}\bigr)^{2}\\
&\le 2\frac1n\sum_{i=1}^{n}\norm{\grad f_i(x_{t}^{s})-\grad f_i(x_{t-1}^{s})}^{2}+2\norm{\grad f(x_{t}^{s})-\grad f(x_{t-1}^{s})}^{2}\\
&\le 2L^{2}\norm{x_{t}^{s}-x_{t-1}^{s}}^{2}+2L^{2}\norm{x_{t}^{s}-x_{t-1}^{s}}^{2}=2L^{2}\norm{x_{t}^{s}-x_{t-1}^{s}}^{2}.
\end{aligned}
\]
\end{proof}

\section*{Main Proof (Step‑by‑Step Derivation)}
We prove Theorem~\ref{thm:linear}.  All expectations are taken with respect to the random choices of $i_t$ conditioned on the past $\mathcal{F}_{t}$.

\noindent\textbf{Step 1. Smoothness descent inequality.}  
For any $x,y$,
\[
f(y)\le f(x)+\inner{\grad f(x)}{y-x}+\frac{L}{2}\norm{y-x}^{2}.
\tag{5}
\]  

\noindent\textbf{Step 2. Apply (5) with $x=x_{t}^{s}$ and $y=x_{t+1}^{s}=x_{t}^{s}-\eta v_{t}^{s}$.}  
\[
\begin{aligned}
f(x_{t+1}^{s})
&\le f(x_{t}^{s})-\eta\inner{\grad f(x_{t}^{s})}{v_{t}^{s}}+\frac{L\eta^{2}}{2}\norm{v_{t}^{s}}^{2}.
\tag{6}
\end{aligned}
\]  

\noindent\textbf{Step 3. Expand the inner product using Lemma~\ref{lem:unbiased}.}  
\[
\inner{\grad f(x_{t}^{s})}{v_{t}^{s}}
= \norm{\grad f(x_{t}^{s})}^{2}
+\inner{\grad f(x_{t}^{s})}{v_{t}^{s}-\grad f(x_{t}^{s})}.
\tag{7}
\]  

\noindent\textbf{Step 4. Take conditional expectation of (6).}  
\[
\begin{aligned}
\E_{i_t}[f(x_{t+1}^{s})\mid\mathcal{F}_{t}]
&\le f(x_{t}^{s})-\eta\norm{\grad f(x_{t}^{s})}^{2}
-\eta\inner{\grad f(x_{t}^{s})}{\E_{i_t}[v_{t}^{s}-\grad f(x_{t}^{s})\mid\mathcal{F}_{t}]}\\
&\qquad+\frac{L\eta^{2}}{2}\E_{i_t}\!\bigl[\norm{v_{t}^{s}}^{2}\mid\mathcal{F}_{t}\bigr].
\end{aligned}
\tag{8}
\]  
By Lemma~\ref{lem:unbiased} the second inner product term vanishes.

\noindent\textbf{Step 5. Bound $\E[\norm{v_{t}^{s}}^{2}]$ via variance decomposition.}  
\[
\begin{aligned}
\E_{i_t}\!\bigl[\norm{v_{t}^{s}}^{2}\mid\mathcal{F}_{t}\bigr]
&= \norm{\grad f(x_{t}^{s})}^{2}
+\E_{i_t}\!\bigl[\norm{v_{t}^{s}-\grad f(x_{t}^{s})}^{2}\mid\mathcal{F}_{t}\bigr] \\
&\stackrel{(a)}{\le} \norm{\grad f(x_{t}^{s})}^{2}+2L^{2}\norm{x_{t}^{s}-x_{t-1}^{s}}^{2},
\end{aligned}
\tag{9}
\]  
where (a) follows from Lemma~\ref{lem:var}.

\noindent\textbf{Step 6. Substitute (9) into (8).}  
\[
\begin{aligned}
\E_{i_t}[f(x_{t+1}^{s})\mid\mathcal{F}_{t}]
&\le f(x_{t}^{s})-\eta\norm{\grad f(x_{t}^{s})}^{2}
+\frac{L\eta^{2}}{2}\norm{\grad f(x_{t}^{s})}^{2}
+\frac{L\eta^{2}}{2}\,2L^{2}\norm{x_{t}^{s}-x_{t-1}^{s}}^{2}\\
&= f(x_{t}^{s})-\underbrace{\bigl(\eta-\tfrac{L\eta^{2}}{2}\bigr)}_{\displaystyle \alpha}\norm{\grad f(x_{t}^{s})}^{2}
+L^{3}\eta^{2}\norm{x_{t}^{s}-x_{t-1}^{s}}^{2}.
\end{aligned}
\tag{10}
\]  

\noindent\textbf{Step 7. Relate $\norm{x_{t}^{s}-x_{t-1}^{s}}$ to $v_{t-1}^{s}$.}  
From the update rule (2),
\[
x_{t}^{s}-x_{t-1}^{s}= -\eta v_{t-1}^{s}.
\tag{11}
\]  

\noindent\textbf{Step 8. Insert (11) into (10) and take full expectation.}  
\[
\begin{aligned}
\E[f(x_{t+1}^{s})]
&\le \E[f(x_{t}^{s})]-\alpha\,\E\!\bigl[\norm{\grad f(x_{t}^{s})}^{2}\bigr]
+L^{3}\eta^{3}\E\!\bigl[\norm{v_{t-1}^{s}}^{2}\bigr].
\end{aligned}
\tag{12}
\]  

\noindent\textbf{Step 9. Upper bound $\E[\norm{v_{t-1}^{s}}^{2}]$ using (9) again (now with $t\!\leftarrow\!t-1$).}  
\[
\E\!\bigl[\norm{v_{t-1}^{s}}^{2}\bigr]\le \E\!\bigl[\norm{\grad f(x_{t-1}^{s})}^{2}\bigr]+2L^{2}\E\!\bigl[\norm{x_{t-1}^{s}-x_{t-2}^{s}}^{2}\bigr].
\tag{13}
\]  

\noindent\textbf{Step 10. Recursively substitute (13) into (12) for all inner steps $t=0,\dots,m-1$.}  
After telescoping and noting that $x_{-1}^{s}=\tilde{x}^{s}$, we obtain
\[
\E\!\bigl[f(x_{m}^{s})\bigr]\le \E\!\bigl[f(\tilde{x}^{s})\bigr]
-\bigl(\alpha- L^{3}\eta^{3}\bigr)\sum_{t=0}^{m-1}\E\!\bigl[\norm{\grad f(x_{t}^{s})}^{2}\bigr]
+ C\,\eta^{3}\sum_{t=0}^{m-2}\E\!\bigl[\norm{\grad f(x_{t}^{s})}^{2}\bigr],
\tag{14}
\]  
where $C=2L^{5}$ collects the constants from the second term of (13).  

\noindent\textbf{Step 11. Choose $\eta\le 1/(3L)$ so that $\alpha-\!L^{3}\eta^{3}\ge \tfrac{\eta}{2}$.}  
Indeed,
\[
\alpha-\!L^{3}\eta^{3}
 =\eta-\frac{L\eta^{2}}{2}-L^{3}\eta^{3}
 =\eta\Bigl(1-\frac{L\eta}{2}-L^{2}\eta^{2}\Bigr)
 \ge\eta\Bigl(1-\frac13-\frac19\Bigr)=\frac{\eta}{2}.
\tag{15}
\]  

\noindent\textbf{Step 12. Drop the non‑negative sum term and keep only the first term of (14).}  
\[
\E\!\bigl[f(\tilde{x}^{s+1})\bigr]
\le \E\!\bigl[f(\tilde{x}^{s})\bigr]-\frac{\eta}{2}\sum_{t=0}^{m-1}\E\!\bigl[\norm{\grad f(x_{t}^{s})}^{2}\bigr].
\tag{16}
\]  

\noindent\textbf{Step 13. Apply the PL inequality (Assumption~\ref{ass:pl}) to each $x_{t}^{s}$.}  
\[
\norm{\grad f(x_{t}^{s})}^{2}\ge 2\mu\bigl(f(x_{t}^{s})-f^{\star}\bigr).
\tag{17}
\]  

\noindent\textbf{Step 14. Substitute (17) into (16) and use that $f(x_{t}^{s})\ge f(\tilde{x}^{s+1})$ for all $t$ (monotonicity of the inner loop follows from (6) with $\eta\le 1/L$).}  
\[
\begin{aligned}
\E\!\bigl[f(\tilde{x}^{s+1})\bigr]
&\le \E\!\bigl[f(\tilde{x}^{s})\bigr]-\eta\mu\sum_{t=0}^{m-1}\E\!\bigl[f(x_{t}^{s})-f^{\star}\bigr]\\
&\le \E\!\bigl[f(\tilde{x}^{s})\bigr]-\eta\mu\,m\,\E\!\bigl[f(\tilde{x}^{s+1})-f^{\star}\bigr].
\end{aligned}
\tag{18}
\]  

\noindent\textbf{Step 15. Rearrange (18) to isolate $\E[f(\tilde{x}^{s+1})-f^{\star}]$.}  
\[
\E\!\bigl[f(\tilde{x}^{s+1})-f^{\star}\bigr]
\le \bigl(1-\eta\mu\bigr)^{m}\,\E\!\bigl[f(\tilde{x}^{s})-f^{\star}\bigr].
\tag{19}
\]  
Equation (19) is exactly the recursion (3) claimed in Theorem~\ref{thm:linear}. Repeatedly applying it yields the final bound after $S$ outer epochs. ∎

\section*{Rate Derivation (With Constants)}
From (19) let $\rho:=1-\eta\mu\in(0,1)$. After $S$ outer epochs,
\[
\E\!\bigl[f(\tilde{x}^{S})-f^{\star}\bigr]\le \rho^{mS}\bigl[f(x^{0})-f^{\star}\bigr].
\]
Choosing the maximal admissible stepsize $\eta=1/(3L)$ gives
\[
\rho = 1-\frac{\mu}{3L},
\qquad
\text{so }\;
\text{contraction factor per inner iteration}
= \bigl(1-\tfrac{\mu}{3L}\bigr)^{1/m}.
\]

\section*{Special Cases}
\begin{itemize}
\item \textbf{Strongly convex case.} If $f$ is $\mu$‑strongly convex, the PL condition holds with the same $\mu$, and the same linear rate (3) applies.
\item \textbf{Exact PL without smoothness.} If smoothness is replaced by a weaker Lipschitz gradient on average, the proof can be adapted by using a larger stepsize bound; the linear factor becomes $1-\eta\mu$ with $\eta\le 1/(2L_{\text{avg}})$.
\item \textbf{Full‑gradient limit.} Setting $m=1$ and $i_t$ deterministic yields $v_{t}^{s}=\grad f(x_{t}^{s})$, and (3) reduces to the classic GD linear rate $f(x_{t+1})-f^{\star}\le (1-\eta\mu)(f(x_{t})-f^{\star})$.
\end{itemize}

\end{document}